\documentclass[12pt, a4paper]{article}

\usepackage{amsmath, amssymb, amsfonts}
\usepackage{float}
\usepackage{graphicx}
\usepackage{listings}
\usepackage{xcolor}
\usepackage[linesnumbered, ruled]{algorithm2e}
\usepackage{subcaption}
\usepackage[hidelinks]{hyperref}

\definecolor{codegreen}{rgb}{0,0.6,0}
\definecolor{codegray}{rgb}{0.5,0.5,0.5}
\definecolor{codepurple}{rgb}{0.58,0,0.82}
\hypersetup{linktoc=all}

\lstset{
    language=Python,
    basicstyle=\ttfamily\small,
    keywordstyle=\color{blue},
    commentstyle=\color{codegreen},
    stringstyle=\color{codepurple},
    breaklines=true,
    numbers=left,
    numberstyle=\tiny\color{codegray},
    frame=single,
    showstringspaces=false
}

\title{AI Masterclass Test 2}
\author{}
\date{}


\begin{document}

\section{Q-function}

$$Q(s, a) \leftarrow Q(s, a) + \alpha [ r + \gamma \max_{a'} Q(s', a') - Q(s, a) ]$$
% \begin{verbatim}
%     Q[state[0], state[1], action] += alpha * (reward + 
%         gamma * np.max(Q[next_state[0], next_state[1], :])
%         - Q[state[0], state[1], action])
% \end{verbatim}
Q-learning is off-policy (learns about a different policy than the one it uses to choose actions during training).
\begin{itemize}
    \item The behaviour policy is $\epsilon$-greedy, so $\epsilon$\% of the time it takes a random action. This is a deliberately suboptimal, exploratory policy.
    \item The target policy is full greedy (\texttt{np.max}) and takes the best possible action from $s'$ regardless if the action to get there is random.
\end{itemize}
The Q values converge to the optimal Q values - the values corresponding to always acting greedily - even though the agent spent training acting suboptimally due to exploration.

\end{document}
